problem:  
**Conjecture (Global Hypoelliptic Rigidity Conjecture).**  
Let \( (M,g) \) be a compact smooth manifold equipped with smooth vector fields \( \{X_1,\dots,X_k\} \) satisfying Hörmander’s bracket‑generating condition, and let  
\[
\Delta_H = -\sum_{i=1}^k X_i^*X_i
\]
be the associated sub‑Laplacian.  
Assume \( \Delta_H \) is globally hypoelliptic on \(M\) but not elliptic.

*Conjecture:* Such a non‑elliptic but globally hypoelliptic Hörmander system can occur **only** when \(M\) is a compact homogeneous manifold \(G/\Gamma\) (e.g. a torus or a compact quotient of a Lie group by a lattice), and the vector fields \(X_i\) arise from a transitive Lie algebra action of \(G\).  

In particular, global hypoellipticity is conjectured to force the manifold to admit a transitive symmetry group compatible with the Hörmander frame.  

*Refinement:* Investigate whether global hypoellipticity implies the existence of a representation‑theoretic spectral decomposition for \(\Delta_H\) corresponding to unitary representations of \(G\), thereby establishing rigidity of both geometric and analytic types.

Why is it a "good" problem:  
This conjecture proposes a precise rigidity principle connecting **microlocal analysis**, **Lie theory**, and **global differential geometry**. It refines the spirit of the Greenfield–Wallach conjecture from single vector fields to Hörmander sub‑Laplacians, extending the analytic phenomenon of global hypoellipticity into the subelliptic realm. Although simple to state, resolving it would require new deep insights into how global analytic regularity enforces algebraic or geometric symmetry. It would bridge the study of hypoelliptic operators, spectral geometry on non‑elliptic structures, and the classification of compact manifolds by symmetry. The problem stands at a point where known examples (tori, homogeneous spaces) fit naturally, but no general theorem confirms the exclusivity, making it both conceptually clear and genuinely open—hence, an excellent “good” problem.